\documentclass[11pt]{standalone}
\usepackage{amssymb}
\usepackage{amsmath}
\usepackage[no-math]{fontspec}
\usepackage{unicode-math}
\setmainfont{Libertinus Sans}
\setsansfont{Libertinus Sans}
\setmathfont{Libertinus Math}
\usepackage{pgf,xcolor}
\usepackage{tikz}
\usetikzlibrary{arrows.meta, backgrounds, calc}
\usepackage{pgfplots}
\pgfplotsset{compat=newest}
\usepackage{pifont}

\begin{document}

\definecolor{my_darkgray}{HTML}{484949}
\definecolor{my_lightgray}{HTML}{F3F4F4}
\definecolor{my_darkblue}{HTML}{005A94}
\definecolor{my_lightblue}{HTML}{CCDEE9}
\definecolor{my_yellow}{HTML}{FFEC7F}
\definecolor{my_red}{HTML}{E60000}
\definecolor{my_orange}{HTML}{E87846}

\begin{tikzpicture}[
  show background rectangle,
  background rectangle/.style={fill=my_lightgray, rounded corners=8pt},
  inner frame sep=0.8cm,
  annotation/.style={draw=my_darkgray, thick, -{Stealth[length=7pt, width=5pt]}}
]

%-----------------------------------------------------------------------
% Knotenkoordinaten (Rautenform)
%-----------------------------------------------------------------------
\coordinate (A) at (4, 6);    % oberer Knoten
\coordinate (B) at (1, 3);    % Knoten Mitte links
\coordinate (C) at (7, 3);    % Knoten Mitte rechts
\coordinate (D) at (4, 0);    % unterer Knoten

%-----------------------------------------------------------------------
% Leitungen (zuerst, damit Widerstandsboxen darübergezeichnet werden)
%-----------------------------------------------------------------------
\draw[draw=my_darkblue, ultra thick] (A) -- (B);   % linker Zweig oben
\draw[draw=my_darkblue, ultra thick] (B) -- (D);   % linker Zweig unten
\draw[draw=my_darkblue, ultra thick] (A) -- (C);   % rechter Zweig oben
\draw[draw=my_darkblue, ultra thick] (C) -- (D);   % rechter Zweig unten
\draw[draw=my_darkblue, ultra thick] (B) -- (C);   % Brückenzweig
% Außenleitung mit Spannungsquelle (rechts) – mit Lücke für Batteriesymbol:
\draw[draw=my_darkblue, ultra thick]
    (D) -- (9.0, 0) -- (9.0, 2.8);
\draw[draw=my_darkblue, ultra thick]
    (9.0, 3.2) -- (9.0, 6) -- (A);

%-----------------------------------------------------------------------
% Widerstandsboxen
% Jede Box ist ein my_lightblue-Rechteck, zentriert am Mittelpunkt des
% jeweiligen Drahtstücks, rotiert parallel zur Drahtrichtung.
% Halbbreite (entlang Draht): 0.6, Halbhöhe (quer): 0.28
%-----------------------------------------------------------------------

% R1: Mittelpunkt A(4,6) -> B(1,3) = (2.5, 4.5)
% Drahtrichtung SW (225 deg) => Rotation 45 deg
\filldraw[fill=my_lightblue, draw=my_darkblue, thick,
          rotate around={45:(2.5,4.5)}]
    (1.9,4.22) rectangle (3.1,4.78);
\node[font=\small\bfseries, text=my_darkblue,
      rotate=45] at (2.5, 4.5) {$R_1$};

% R2: Mittelpunkt B(1,3) -> D(4,0) = (2.5, 1.5)
% Drahtrichtung SE (-45 deg) => Rotation -45 deg
\filldraw[fill=my_lightblue, draw=my_darkblue, thick,
          rotate around={-45:(2.5,1.5)}]
    (1.9,1.22) rectangle (3.1,1.78);
\node[font=\small\bfseries, text=my_darkblue,
      rotate=-45] at (2.5, 1.5) {$R_2$};

% R3: Mittelpunkt A(4,6) -> C(7,3) = (5.5, 4.5)
% Drahtrichtung SE (-45 deg) => Rotation -45 deg
\filldraw[fill=my_lightblue, draw=my_darkblue, thick,
          rotate around={-45:(5.5,4.5)}]
    (4.9,4.22) rectangle (6.1,4.78);
\node[font=\small\bfseries, text=my_darkblue,
      rotate=-45] at (5.5, 4.5) {$R_3$};

% R4: Mittelpunkt C(7,3) -> D(4,0) = (5.5, 1.5)
% Drahtrichtung SW (225 deg) => Rotation 45 deg
% Hervorgehoben (my_yellow): R4 ist der variable Messwiderstand
\filldraw[fill=my_yellow, draw=my_darkblue, thick,
          rotate around={45:(5.5,1.5)}]
    (4.9,1.22) rectangle (6.1,1.78);
\node[font=\small\bfseries, text=my_darkblue,
      rotate=45] at (5.5, 1.5) {$R_4$};
% Hinweis: R4 ist variabel (Messwiderstand)
\node[font=\footnotesize, text=my_darkgray, anchor=north west]
    at (5.5, 1.3) {(variabel)};

% RL: Mittelpunkt B(1,3) -> C(7,3) = (4, 3)
% Horizontal => keine Rotation
\filldraw[fill=my_lightblue, draw=my_darkblue, thick]
    (3.3, 2.65) rectangle (4.7, 3.35);
\node[font=\small\bfseries, text=my_darkblue]
    at (4, 3) {$R_B$};

% U0: Batteriesymbol – Striche direkt an den Zuleitungen, kein Spalt
% - Pol (kurz, dick) bei y=2.8, direkt am unteren Drahtende
% + Pol (lang, dünn) bei y=3.2, direkt am oberen Drahtanfang
\draw[draw=my_darkblue, line width=3pt] (8.75, 2.8) -- (9.25, 2.8);   % - Pol
\draw[draw=my_darkblue, line width=3pt]  (8.55, 3.2) -- (9.45, 3.2);   % + Pol
% Beschriftung rechts neben dem Symbol
\node[font=\small\bfseries, text=my_darkblue, anchor=west]
    at (9.55, 3.0) {$U_0$};
\node[font=\small\bfseries, text=my_darkgray, anchor=west]
    at (8.55, 3.5) {$+$};
\node[font=\small\bfseries, text=my_darkgray, anchor=west]
    at (8.55, 2.5) {$-$};

%-----------------------------------------------------------------------
% Knotenlabels A, B, C, D
%-----------------------------------------------------------------------
\node[font=\small\bfseries, text=my_darkgray, anchor=south]
    at (4, 6.18) {$K_1$};
\node[font=\small\bfseries, text=my_darkgray, anchor=east]
    at (0.82, 3) {$K_2$};
\node[font=\small\bfseries, text=my_darkgray, anchor=west]
    at (7.18, 3) {$K_3$};
%\node[font=\small\bfseries, text=my_darkgray, anchor=north]
%    at (4, -0.18) {$D$};

%-----------------------------------------------------------------------
% Knotenpunkte (farbige Kreise)
%-----------------------------------------------------------------------
\foreach \p in {A, B, C, D} {
    \filldraw[fill=my_orange, draw=my_lightgray, thick] (\p) circle (0.13);
}

%-----------------------------------------------------------------------
% Stromrichtungspfeile
% Platziert zwischen Knoten und Widerstandsbox (0.15 bis 0.33 der Strecke)
%-----------------------------------------------------------------------

% I0: auf der oberen Außenleitung, kurz vor Knoten A (von rechts kommend)
\draw[annotation] (5.2, 6) -- (4.4, 6);
\node[font=\small, text=my_darkgray, anchor=south] at (4.8, 6.08) {$I_0$};

% I1 (Wiederholung): kurz vor Knoten B, für Knotenregel in B
\draw[annotation] ($(A)!0.67!(B)$) -- ($(A)!0.82!(B)$);
\node[font=\small, text=my_darkgray, anchor=east]
    at ($(A)!0.75!(B)+(-0.15, 0.1)$) {$I_1$};

% I3 (Wiederholung): kurz vor Knoten C, für Knotenregel in C
\draw[annotation] ($(A)!0.67!(C)$) -- ($(A)!0.82!(C)$);
\node[font=\small, text=my_darkgray, anchor=west]
    at ($(A)!0.75!(C)+(0.15, 0.1)$) {$I_3$};

% I4 (Wiederholung): kurz nach Knoten C, für Knotenregel in C
\draw[annotation] ($(C)!0.15!(D)$) -- ($(C)!0.30!(D)$);
\node[font=\small, text=my_darkgray, anchor=west]
    at ($(C)!0.22!(D)+(0.1, 0.0)$) {$I_4$};
\draw[annotation] ($(A)!0.15!(B)$) -- ($(A)!0.33!(B)$);
\node[font=\small, text=my_darkgray, anchor=east]
    at ($(A)!0.24!(B)+(-0.15, 0.1)$) {$I_1$};

% I2: B -> D (linker Zweig unten), vor dem Widerstand
\draw[annotation] ($(B)!0.15!(D)$) -- ($(B)!0.33!(D)$);
\node[font=\small, text=my_darkgray, anchor=east]
    at ($(B)!0.24!(D)+(-0.15, 0.0)$) {$I_2$};

% I3: A -> C (rechter Zweig oben), vor dem Widerstand
\draw[annotation] ($(A)!0.15!(C)$) -- ($(A)!0.33!(C)$);
\node[font=\small, text=my_darkgray, anchor=west]
    at ($(A)!0.24!(C)+(0.15, 0.1)$) {$I_3$};

% I4: C -> D (rechter Zweig unten), nach dem Widerstand
\draw[annotation] ($(C)!0.68!(D)$) -- ($(C)!0.82!(D)$);
\node[font=\small, text=my_darkgray, anchor=west]
    at ($(C)!0.70!(D)+(0.1, -0.3)$) {$I_4$};

% I: B -> C (Brückenzweig), nach dem Widerstand
\draw[annotation] ($(B)!0.63!(C)$) -- ($(B)!0.76!(C)$);
\node[font=\small, text=my_darkgray, anchor=south]
    at ($(B)!0.695!(C)+(0, 0.18)$) {$I$};

\end{tikzpicture}

\end{document}
